\section{AHSME 1970}
        \begin{enumerate}
        	\item (\textit{AHSME 1970}) The fourth power of $\sqrt{1+\sqrt{1+\sqrt{1}}}$ is:


 $\textbf{(A) }\sqrt{2}+\sqrt{3}\qquad \textbf{(B) }\frac{1}{2}(7+3\sqrt{5})\qquad \textbf{(C) }1+2\sqrt{3}\qquad \textbf{(D) }3\qquad \textbf{(E) }3+2\sqrt{2}$


 

	\item (\textit{AHSME 1970}) A square and a circle have equal perimeters. The ratio of the area of the circle to the area of the square is:


 $\textbf{(A) }\frac{4}{\pi}\qquad \textbf{(B) }\frac{\pi}{\sqrt{2}}\qquad \textbf{(C) }\frac{4}{1}\qquad \textbf{(D) }\frac{\sqrt{2}}{\pi}\qquad \textbf{(E) }\frac{\pi}{4}$


 

	\item (\textit{AHSME 1970}) If $x=1+2^p$ and $y=1+2^{-p}$, then $y$ in terms of $x$ is:


 $\textbf{(A) }\dfrac{x+1}{x-1}\qquad \textbf{(B) }\dfrac{x+2}{x-1}\qquad \textbf{(C) }\dfrac{x}{x-1}\qquad \textbf{(D) }2-x\qquad \textbf{(E) }\frac{x-1}{x}$


 

	\item (\textit{AHSME 1970}) Let $S$ be the set of all numbers which are the sum of the squares of three consecutive integers. Then we can say that:


 $\textbf{(A) }\text{No member of }S\text{ is divisible by }2\qquad\\ \textbf{(B) }\text{No member of }S\text{ is divisible by }3\text{ but some member is divisible by }11\qquad\\ \textbf{(C) }\text{No member of }S\text{ is divisible by }3\text{ or }5\qquad\\ \textbf{(D) }\text{No member of }S\text{ is divisible by }3\text{ or }7\qquad\\ \textbf{(E) }\text{None of these}$


 

	\item (\textit{AHSME 1970}) If $f(x)=\dfrac{x^4+x^2}{x+1}$, then $f(i)$, where $i=\sqrt{-1}$, is equal to:


 $\textbf{(A) }1+i\qquad \textbf{(B) }1\qquad \textbf{(C) }-1\qquad \textbf{(D) }0\qquad  \textbf{(E) }-1-i$


 

	\item (\textit{AHSME 1970}) The smallest value of $x^2+8x$ for real values of $x$ is:


 $\textbf{(A) }-16.25\qquad \textbf{(B) }-16\qquad \textbf{(C) }-15\qquad \textbf{(D) }-8\qquad  \textbf{(E) }\text{None of these}$


 

	\item (\textit{AHSME 1970}) Inside square $ABCD$ with side $s$, quarter-circle arcs with radii $s$ and centers at $A$ and $B$ are drawn. 
These arcs intersect at point $X$ inside the square. How far is $X$ from side $CD$?


 $\textbf{(A) }\frac{1}{2}s(\sqrt{3}+4)\qquad \textbf{(B) }\frac{1}{2}s\sqrt{3}\qquad \textbf{(C) }\frac{1}{2}s(1+\sqrt{3})\qquad\\ \textbf{(D) }\frac{1}{2}s(\sqrt{3}-1)\qquad  \textbf{(E) }\frac{1}{2}s(2-\sqrt{3})$


 

	\item (\textit{AHSME 1970}) If $a=\log_8{225}$ and $b=\log_2{15}$, then


 $\textbf{(A) }a=\frac{1}{2}b\qquad \textbf{(B) }a=\frac{2b}{3}\qquad \textbf{(C) }a=b\qquad \textbf{(D) }b=\frac{1}{2}a\qquad \textbf{(E) }a=\frac{3b}{2}$


 

	\item (\textit{AHSME 1970}) Points $P$ and $Q$ are on line segment $AB$, and both points are on the same side of the midpoint of $AB$.
Point $P$ divides $AB$ in the ratio $2:3$ and $Q$ divides $AB$ in the ratio $3:4$. If $PQ=2$, then the length of segment $AB$ is


 $\textbf{(A) }12\qquad \textbf{(B) }28\qquad \textbf{(C) }70\qquad \textbf{(D) }75\qquad  \textbf{(E) }105$


 

	\item (\textit{AHSME 1970}) Let $F=.48181\cdots$ be an infinite repeating decimal with the digits $8$ and $1$ repeating. 
When $F$ is written as a fraction in lowest terms, the denominator exceeds the numerator by


 $\textbf{(A) }13\qquad \textbf{(B) }14\qquad \textbf{(C) }29\qquad \textbf{(D) }57\qquad  \textbf{(E) }126$


 

	\item (\textit{AHSME 1970}) If two factors of $2x^3-hx+k$ are $x+2$ and $x-1$, the value of $|2h-3k|$ is


 $\textbf{(A) }4\qquad \textbf{(B) }3\qquad \textbf{(C) }2\qquad \textbf{(D) }1\qquad  \textbf{(E) }0$


 

	\item (\textit{AHSME 1970}) A circle with radius $r$ is tangent to sides $AB$, $AD$, and $CD$ of rectangle $ABCD$ and passes through
the midpoint of diagonal $AC$. The area of the rectangle in terms of $r$, is


 $\textbf{(A) }4r^2\qquad \textbf{(B) }6r^2\qquad \textbf{(C) }8r^2\qquad \textbf{(D) }12r^2\qquad  \textbf{(E) }20r^2$


 

	\item (\textit{AHSME 1970}) Given the binary operation $\ast$ defined by $a\ast b=a^b$ for all positive numbers $a$ and $b$. The for all positive $a,b,c,n$, we have


 $\textbf{(A) }a\ast b=b\ast a\qquad \textbf{(B) }a\ast (b\ast c)=(a\ast b)\ast c\qquad\\ \textbf{(C) }(a\ast b^n)=(a\ast n)\ast b\qquad \textbf{(D) }(a\ast b)^n=a\ast (bn)\qquad  \textbf{(E) }\text{None of these}$


 

	\item (\textit{AHSME 1970}) Consider $x^2+px+q=0$ where $p$ and $q$ are positive numbers. If the roots of this equation differ by $1$, then $p$ equals


 $\textbf{(A) }\sqrt{4q+1}\qquad \textbf{(B) }q-1\qquad \textbf{(C) }-\sqrt{4q+1}\qquad\\ \textbf{(D) }q+1\qquad \textbf{(E) }\sqrt{4q-1}$


 

	\item (\textit{AHSME 1970}) Lines in the $xy$-plane are drawn through the point $(3,4)$ and the trisection points of the line segment 
joining the points $(-4,5)$ and $(5,-1)$. One of these lines has the equation 


 $\textbf{(A) }3x-2y-1=0\qquad \textbf{(B) }4x-5y+8=0\qquad \textbf{(C) }5x+2y-23=0\qquad\\ \textbf{(D) }x+7y-31=0\qquad  \textbf{(E) }x-4y+13=0$


 

	\item (\textit{AHSME 1970}) If $F(n)$ is a function such that $F(1)=F(2)=F(3)=1$, and such that $F(n+1)=\dfrac{F(n)\cdot F(n-1)+1}{F(n-2)}$ for $n\ge 3$,
then $F(6)$ is equal to


 $\textbf{(A) }2\qquad \textbf{(B) }3\qquad \textbf{(C) }7\qquad \textbf{(D) }11\qquad  \textbf{(E) }26$


 

	\item (\textit{AHSME 1970}) If $r\ge 0$, then for all $p$ and $q$ such that $pq\neq 0$ and $pr>qr$, we have


 $\textbf{(A) }-p>-q\qquad \textbf{(B) }-p>q\qquad \textbf{(C) }1>-q/p\qquad\\ \textbf{(D) }1<q/p\qquad \textbf{(E) }\text{None of These}$


 

	\item (\textit{AHSME 1970}) $\sqrt{3+2\sqrt{2}}-\sqrt{3-2\sqrt{2}}$ is equal to


 $\textbf{(A) }2\qquad \textbf{(B) }2\sqrt{3}\qquad \textbf{(C) }4\sqrt{2}\qquad \textbf{(D) }\sqrt{6}\qquad  \textbf{(E) }2\sqrt{2}$


 

	\item (\textit{AHSME 1970}) The sum of an infinite geometric series with common ratio $r$ such that $|r|<1$, is $15$,
and the sum of the squares of the terms of this series is $45$. The first term of the series is


 $\textbf{(A) }12\qquad \textbf{(B) }10\qquad \textbf{(C) }5\qquad \textbf{(D) }3\qquad  \textbf{(E) }2$


 

	\item (\textit{AHSME 1970}) Lines $HK$ and $BC$ lie in a plane. $M$ is the midpoint of line segment $BC$, and $BH$ and $CK$ are perpendicular to $HK$. Then we 


 $\textbf{(A) }\text{always have }MH=MK\qquad\\ \textbf{(B) }\text{always have }MH>BK\qquad\\ \textbf{(C) }\text{sometimes have }MH=MK\text{ but not always}\qquad\\ \textbf{(D) }\text{always have }MH>MB\qquad \\ \textbf{(E) }\text{always have }BH<BC$


 

	\item (\textit{AHSME 1970}) On an auto trip, the distance read from the instrument panel was $450$ miles. 
With snow tires on for the return trip over the same route, the reading was $440$ miles. 
Find, to the nearest hundredth of an inch, the increase in radius of the wheels if the original radius was $15$ inches.


 $\textbf{(A) }.33\qquad \textbf{(B) }.34\qquad \textbf{(C) }.35\qquad \textbf{(D) }.38\qquad  \textbf{(E) }.66$


 

	\item (\textit{AHSME 1970}) If the sum of the first $3n$ positive integers is $150$ more than the sum of the first $n$ positive integers,
then the sum of the first $4n$ positive integers is


 $\textbf{(A) }300\qquad \textbf{(B) }350\qquad \textbf{(C) }400\qquad \textbf{(D) }450\qquad  \textbf{(E) }600$


 

	\item (\textit{AHSME 1970}) The number $10!$ ($10$ is written in base $10$), when written in the base $12$ system, ends in exactly $k$ zeroes. The value of $k$ is


 $\textbf{(A) }1\qquad \textbf{(B) }2\qquad \textbf{(C) }3\qquad \textbf{(D) }4\qquad  \textbf{(E) } 5$


 

	\item (\textit{AHSME 1970}) An equilateral triangle and a regular hexagon have equal perimeters. If the area of the triangle is $2$, then the area of the hexagon is


 $\textbf{(A) }2\qquad \textbf{(B) }3\qquad \textbf{(C) }4\qquad \textbf{(D) }6\qquad  \textbf{(E) }12$


 

	\item (\textit{AHSME 1970}) For every real number $x$, let $[x]$ be the greatest integer less than or equal to $x$. 
If the postal rate for first class mail is six cents for every ounce or portion thereof, 
then the cost in cents of first-class postage on a letter weighing $W$ ounces is always


 $\textbf{(A) }6W\qquad \textbf{(B) }6[W]\qquad \textbf{(C) }6([W]-1)\qquad \textbf{(D) }6([W]+1)\qquad  \textbf{(E) }-6[-W]$


 

	\item (\textit{AHSME 1970}) The number of distinct points in the $xy$-plane common to the graphs of $(x+y-5)(2x-3y+5)=0$ and $(x-y+1)(3x+2y-12)=0$ is


 $\textbf{(A) }0\qquad \textbf{(B) }1\qquad \textbf{(C) }2\qquad \textbf{(D) }3\qquad  \textbf{(E) }4$


 

	\item (\textit{AHSME 1970}) In a triangle, the area is numerically equal to the perimeter. What is the radius of the inscribed circle?


 $\textbf{(A) }2\qquad \textbf{(B) }3\qquad \textbf{(C) }4\qquad \textbf{(D) }5\qquad  \textbf{(E) }6$


 

	\item (\textit{AHSME 1970}) In triangle $ABC$, the median from vertex $A$ is perpendicular to the median from vertex $B$. 
If the lengths of sides $AC$ and $BC$ are $6$ and $7$ respectively, then the length of side $AB$ is


 $\textbf{(A) }\sqrt{17}\qquad \textbf{(B) }4\qquad \textbf{(C) }4\dfrac{1}{2}\qquad \textbf{(D) }2\sqrt{5}\qquad \textbf{(E) }4\frac{1}{4}$


 

	\item (\textit{AHSME 1970}) It is now between $10:00$ and $11:00$ o'clock, and six minutes from now, the minute hand of the 
watch will be exactly opposite the place where the hour hand was three minutes ago. What is the exact time now?


 $\textbf{(A) }10:05\dfrac{5}{11}\qquad \textbf{(B) }10:07\dfrac{1}{2}\qquad \textbf{(C) }10:10\qquad\\ \textbf{(D) }10:15\qquad \textbf{(E) }10:17\dfrac{1}{2}$


 

	\item (\textit{AHSME 1970}) 
                    \begin{figure}[H]
                        \centering
                        \includegraphics[width=0.5\textwidth]{Problem/AHSME_I_1970_30_Problem_Figure_1.png}
                    \end{figure}
                    


 In the accompanying figure, segments $AB$ and $CD$ are parallel, 
the measure of $\angle{D}$ is twice the measure of $\angle{B}$, and the measures of segments $AD$ and $CD$
are $a$ and $b$ respectively. Then the measure of $AB$ is equal to


 $\textbf{(A) }\dfrac{1}{2}a+2b\qquad \textbf{(B) }\dfrac{3}{2}b+\dfrac{3}{4}a\qquad \textbf{(C) }2a-b\qquad \textbf{(D) }4b-\frac{1}{2}a\qquad \textbf{(E) }a+b$


 

	\item (\textit{AHSME 1970}) If a number is selected at random from the set of all five-digit numbers in which the sum of the digits is equal to $43$,
what is the probability that this number is divisible by $11$?


 $\textbf{(A) }2/5\qquad \textbf{(B) }1/5\qquad \textbf{(C) }1/6\qquad \textbf{(D) }1/11\qquad  \textbf{(E) }1/15$


 

	\item (\textit{AHSME 1970}) $A$ and $B$ travel around a circular track at uniform speeds in opposite directions, 
starting from diametrically opposite points. If they start at the same time, 
meet first after $B$ has travelled $100$ yards, and meet a second time $60$ yards before $A$ completes one lap, 
then the circumference of the track in yards is


 $\textbf{(A) }400\qquad \textbf{(B) }440\qquad \textbf{(C) }480\qquad \textbf{(D) }560\qquad  \textbf{(E) }880$


 

	\item (\textit{AHSME 1970}) Find the sum of the digits of all numerals in the sequence $1,2,3,4,\cdots ,10000$.


 $\textbf{(A) }180,001\qquad \textbf{(B) }154,756\qquad \textbf{(C) }45,001\qquad \textbf{(D) }154,755\qquad  \textbf{(E) }270,001$


 

	\item (\textit{AHSME 1970}) The greatest integer that will divide $13511, 13903$, and $14589$ and leave the same remainder is


 $\textbf{(A) }28\qquad \textbf{(B) }49\qquad \textbf{(C) }98\qquad\\ \textbf{(D) }\text{an odd multiple of }7\text{ greater than }49\qquad\\ \textbf{(E) }\text{an even multiple of }7\text{ greater than }98$


 

	\item (\textit{AHSME 1970}) A retiring employee receives an annual pension proportional to the square root of the number of years of his service. 
Had he served $a$ years more, his pension would have been $p$ dollars greater, whereas, had he served $b$ years more $b\neq a$, 
his pension would have been $q$ dollars greater than the original annual pension. Find his annual pension in terms of $a,b,p$, and $q$.


 $\textbf{(A) }\dfrac{p^2-q^2}{2(a-b)}\qquad \textbf{(B) }\dfrac{(p-q)^2}{2\sqrt{ab}}\qquad \textbf{(C) }\frac{ap^2-bq^2}{2(ap-bq)}\qquad \textbf{(D) }\frac{aq^2-bp^2}{2(bp-aq)}\qquad \textbf{(E) }\sqrt{(a-b)(p-q)}$


 

\end{enumerate}